%%%%%%%%%%%%%%%%%%%%%%%%%%%%%%%%%%%%%%%%%%%%%%%%%%%%%%%%%%%%%%%%%%%%%%%%%%%%%%%
%% TKS v7.4 HIGHER CATEGORIES TRACK
%% $(\infty,1)$-Categories, Quasi-Categories, and Noetic Homotopy Theory
%% Part of TKS v7.4 - December 2025
%% ADDITIVE to v7.3 canon
%%%%%%%%%%%%%%%%%%%%%%%%%%%%%%%%%%%%%%%%%%%%%%%%%%%%%%%%%%%%%%%%%%%%%%%%%%%%%%%

%% ============================================================================
%% CHAPTER 1: QUASI-CATEGORIES AND NOETIC SPACES
%% ============================================================================

\chapter{Quasi-Categories and Noetic Spaces}
\label{ch:quasi-categories}

\begin{tcolorbox}[colback=blue!5!white,colframe=blue!75!black,title=\vnewfour]
This chapter introduces quasi-categories as a model for $(\infty,1)$-categories
and constructs the noetic quasi-category $\mathcal{N}_\infty$.
\end{tcolorbox}

\section{Simplicial Preliminaries}

\begin{definition}[Simplex Category]
\label{def:simplex-cat}
The \textbf{simplex category} $\Delta$ has:
\begin{itemize}
  \item Objects: Finite ordinals $[n] = \{0, 1, \ldots, n\}$ for $n \geq 0$
  \item Morphisms: Order-preserving maps
\end{itemize}
\end{definition}

\begin{definition}[Simplicial Set]
\label{def:simplicial-set}
A \textbf{simplicial set} is a functor $X : \Delta^{\mathrm{op}} \to \mathbf{Set}$.
We write $X_n := X([n])$ for the set of $n$-simplices.
\end{definition}

\begin{definition}[Horn]
\label{def:horn}
For $0 \leq k \leq n$, the \textbf{$k$-horn} $\Lambda^n_k \subset \Delta^n$ is the
simplicial subset generated by all faces except the $k$-th face.
\end{definition}

\section{Definition of Quasi-Categories}

\begin{definition}[Quasi-Category]
\label{def:quasi-category}
A simplicial set $\mathcal{C}$ is a \textbf{quasi-category} (or $\infty$-category)
if it has the \emph{inner horn filling property}: for all $0 < k < n$, every map
$\Lambda^n_k \to \mathcal{C}$ extends to $\Delta^n \to \mathcal{C}$.
\[
\begin{tikzcd}
\Lambda^n_k \arrow[r] \arrow[d, hook] & \mathcal{C} \\
\Delta^n \arrow[ur, dashed] &
\end{tikzcd}
\]
\end{definition}

\begin{remark}[Interpretation]
In a quasi-category:
\begin{itemize}
  \item 0-simplices are objects
  \item 1-simplices are morphisms
  \item 2-simplices witness composition (up to homotopy)
  \item Higher simplices encode coherence data
\end{itemize}
\end{remark}

\section{The Noetic Quasi-Category $\mathcal{N}_\infty$}

\begin{construction}[Noetic Quasi-Category]
\label{constr:noetic-quasi-cat}
The \textbf{noetic quasi-category} $\mathcal{N}_\infty$ is defined by:
\begin{itemize}
  \item $(\mathcal{N}_\infty)_0 = \Ideas$ (Ideas are objects)
  \item $(\mathcal{N}_\infty)_1 = \{(\iota, \iota', \Phi) : \Phi \text{ is a noetic path from } \iota \text{ to } \iota'\}$
  \item $(\mathcal{N}_\infty)_n$ consists of $n$-chains of noetic paths with coherent homotopies
\end{itemize}
\end{construction}

\begin{theorem}[Noetic Quasi-Category is Well-Defined]
\label{thm:noetic-qcat-well-defined}
$\mathcal{N}_\infty$ is a quasi-category. Inner horn fillings correspond to
fractal interpolation between noetic paths.
\end{theorem}

\section{Mapping Spaces and Hom-Objects}

\begin{definition}[Mapping Space]
\label{def:mapping-space}
For $\iota, \iota' \in \mathcal{N}_\infty$, the \textbf{mapping space} is:
\[
\mathrm{Map}_{\mathcal{N}_\infty}(\iota, \iota') := \{\sigma \in (\mathcal{N}_\infty)_\bullet : d_0\sigma = \iota', d_n\sigma = \iota\}
\]
This is a Kan complex.
\end{definition}

\begin{proposition}[Noetic Mapping Spaces]
\label{prop:noetic-mapping}
The mapping space $\mathrm{Map}_{\mathcal{N}_\infty}(\iota, \iota')$ has:
\begin{enumerate}[label=(\roman*)]
  \item 0-simplices: noetic morphisms $\iota \to \iota'$
  \item 1-simplices: homotopies between morphisms
  \item $n$-simplices: higher homotopy data
\end{enumerate}
\end{proposition}

\section{Kan Complexes and Fibrancy}

\begin{definition}[Kan Complex]
\label{def:kan-complex}
A simplicial set $K$ is a \textbf{Kan complex} if all horns (inner and outer)
can be filled.
\end{definition}

\begin{proposition}[Mapping Spaces are Kan]
\label{prop:mapping-kan}
For any quasi-category $\mathcal{C}$ and objects $x, y$, the mapping space
$\mathrm{Map}_\mathcal{C}(x, y)$ is a Kan complex.
\end{proposition}

%% ============================================================================
%% CHAPTER 2: SEGAL SPACES FOR NOETICS
%% ============================================================================

\chapter{Segal Spaces for Noetics}
\label{ch:segal-spaces}

\begin{tcolorbox}[colback=blue!5!white,colframe=blue!75!black,title=\vnewfour]
This chapter develops Segal spaces as an alternative model equivalent to
quasi-categories.
\end{tcolorbox}

\section{Segal Spaces}

\begin{definition}[Segal Space]
\label{def:segal-space}
A \textbf{Segal space} is a simplicial space $X : \Delta^{\mathrm{op}} \to \mathbf{Space}$
satisfying the \emph{Segal condition}: the map
\[
X_n \to X_1 \times_{X_0} X_1 \times_{X_0} \cdots \times_{X_0} X_1
\]
is a weak equivalence for all $n \geq 2$.
\end{definition}

\begin{definition}[Complete Segal Space]
\label{def:complete-segal}
A Segal space $X$ is \textbf{complete} (or \textbf{Rezk complete}) if the map
\[
X_0 \to X_1^{\mathrm{eq}}
\]
is a weak equivalence, where $X_1^{\mathrm{eq}}$ is the space of equivalences.
\end{definition}

\section{Noetic Segal Space}

\begin{construction}[Noetic Segal Space]
\label{constr:noetic-segal}
The \textbf{noetic Segal space} $\mathcal{N}^{\mathrm{Seg}}$ has:
\begin{itemize}
  \item $\mathcal{N}^{\mathrm{Seg}}_0 = |\Ideas|$ (geometric realization of Idea space)
  \item $\mathcal{N}^{\mathrm{Seg}}_1 = $ space of noetic paths
  \item $\mathcal{N}^{\mathrm{Seg}}_n = $ space of composable $n$-tuples
\end{itemize}
\end{construction}

\begin{theorem}[Noetic Segal Condition]
\label{thm:noetic-segal-condition}
The noetic Segal space satisfies the Segal condition via fractal composition.
\end{theorem}

\section{Equivalence with Quasi-Categories}

\begin{theorem}[Segal-Quasi-Category Equivalence]
\label{thm:segal-qcat-equiv}
There is a Quillen equivalence between:
\begin{itemize}
  \item Complete Segal spaces
  \item Quasi-categories
\end{itemize}
Under this equivalence, $\mathcal{N}^{\mathrm{Seg}} \simeq \mathcal{N}_\infty$.
\end{theorem}

%% ============================================================================
%% CHAPTER 3: $(\infty,1)$-CATEGORIES OF IDEAS
%% ============================================================================

\chapter{$(\infty,1)$-Categories of Ideas}
\label{ch:infty-ideas}

\begin{tcolorbox}[colback=blue!5!white,colframe=blue!75!black,title=\vnewfour]
This chapter develops the full $(\infty,1)$-category $\mathbf{Idea}_\infty$.
\end{tcolorbox}

\section{The $(\infty,1)$-Category $\mathbf{Idea}_\infty$}

\begin{definition}[$(\infty,1)$-Category of Ideas]
\label{def:idea-infty}
The \textbf{$(\infty,1)$-category of Ideas} $\mathbf{Idea}_\infty$ extends the
v6.1 category $\Noetica$ with:
\begin{itemize}
  \item Objects: Ideas $\iota \in \Ideas$
  \item 1-morphisms: Noetic transformations
  \item 2-morphisms: Homotopies between transformations
  \item $n$-morphisms: Higher homotopy coherences
\end{itemize}
\end{definition}

\begin{proposition}[World Filtration]
\label{prop:world-filtration}
$\mathbf{Idea}_\infty$ admits a filtration by Worlds:
\[
\mathbf{Idea}_\infty^{(A)} \hookrightarrow \mathbf{Idea}_\infty^{(B)} \hookrightarrow
\mathbf{Idea}_\infty^{(C)} \hookrightarrow \mathbf{Idea}_\infty^{(D)} = \mathbf{Idea}_\infty
\]
\end{proposition}

\section{Higher Morphisms Between Ideas}

\begin{definition}[Higher Noetic Morphism]
\label{def:higher-morphism}
An \textbf{$n$-morphism} in $\mathbf{Idea}_\infty$ is a map
\[
\alpha : \Delta^n \to \mathbf{Idea}_\infty
\]
with boundary conditions specifying source and target $(n-1)$-morphisms.
\end{definition}

\begin{proposition}[Truncation Recovers $\Noetica$]
\label{prop:truncation}
The 1-truncation $\tau_{\leq 1}\mathbf{Idea}_\infty$ is equivalent to the ordinary
category $\Noetica$ from v6.1.
\end{proposition}

\section{Adjunctions in the $\infty$-Setting}

\begin{definition}[$\infty$-Adjunction]
\label{def:infty-adjunction}
An \textbf{$\infty$-adjunction} between $(\infty,1)$-categories $\mathcal{C}$ and $\mathcal{D}$
consists of functors $F : \mathcal{C} \rightleftarrows \mathcal{D} : G$ with a natural
equivalence of mapping spaces:
\[
\mathrm{Map}_\mathcal{D}(F(x), y) \simeq \mathrm{Map}_\mathcal{C}(x, G(y))
\]
\end{definition}

\begin{theorem}[Noetic Adjunctions]
\label{thm:noetic-adjunctions}
The following are $\infty$-adjunctions on $\mathbf{Idea}_\infty$:
\begin{enumerate}[label=(\roman*)]
  \item \textbf{ACBE adjunction}: World descent $\dashv$ World ascent
  \item \textbf{Fractal adjunction}: Fractal application $\dashv$ Fractal extraction
  \item \textbf{RPM adjunction}: Free RPM monad $\dashv$ Forgetful
\end{enumerate}
\end{theorem}

%% ============================================================================
%% CHAPTER 4: $(\infty,1)$-TOPOI AND NOETIC SHEAVES
%% ============================================================================

\chapter{$(\infty,1)$-Topoi and Noetic Sheaves}
\label{ch:infty-topoi}

\begin{tcolorbox}[colback=blue!5!white,colframe=blue!75!black,title=\vnewfour]
This chapter develops the noetic $\infty$-topos $\mathbf{Noe}_\infty$.
\end{tcolorbox}

\section{$(\infty,1)$-Topoi}

\begin{definition}[Presentable $(\infty,1)$-Category]
\label{def:presentable}
An $(\infty,1)$-category is \textbf{presentable} if it is accessible and admits
all small colimits.
\end{definition}

\begin{definition}[$(\infty,1)$-Topos]
\label{def:infty-topos}
An \textbf{$(\infty,1)$-topos} is a presentable $(\infty,1)$-category $\mathcal{E}$
that is left exact localization of a presheaf category:
\[
\mathcal{E} \simeq L\mathbf{PSh}(\mathcal{C})
\]
for some small $(\infty,1)$-category $\mathcal{C}$.
\end{definition}

\section{The Noetic $\infty$-Topos}

\begin{definition}[Noetic $\infty$-Topos]
\label{def:noetic-infty-topos}
The \textbf{noetic $\infty$-topos} $\InftyTopos_\nu$ is defined as:
\[
\InftyTopos_\nu := \Sh_\infty(\Ideas, \tau_{\mathrm{noetic}})
\]
the $(\infty,1)$-category of $\infty$-sheaves on the noetic site.
\end{definition}

\begin{theorem}[$\InftyTopos_\nu$ is an $\infty$-Topos]
\label{thm:noe-infty-topos}
$\InftyTopos_\nu$ satisfies the axioms of an $(\infty,1)$-topos.
\end{theorem}

\section{Descent and Higher Sheaf Conditions}

\begin{definition}[$\infty$-Descent]
\label{def:infty-descent}
A presheaf $F : \mathcal{C}^{\mathrm{op}} \to \mathbf{Space}$ satisfies \textbf{$\infty$-descent}
for a covering sieve $S$ if the natural map
\[
F(U) \to \lim_{\Delta} F(S_\bullet)
\]
is an equivalence, where $S_\bullet$ is the \v{C}ech nerve.
\end{definition}

\begin{theorem}[Noetic Descent]
\label{thm:noetic-descent}
The noetic topology satisfies effective $\infty$-descent: an $\infty$-presheaf
is a sheaf iff it satisfies descent for all noetic coverings.
\end{theorem}

\section{$\infty$-Giraud Axioms}

\begin{theorem}[$\infty$-Giraud Axioms]
\label{thm:infty-giraud}
An $(\infty,1)$-category $\mathcal{E}$ is an $\infty$-topos iff:
\begin{enumerate}[label=(G\arabic*)]
  \item $\mathcal{E}$ has all small colimits
  \item Colimits are universal
  \item Coproducts are disjoint
  \item Every groupoid object is effective
  \item There exists a set of generators
\end{enumerate}
\end{theorem}

\begin{corollary}[$\InftyTopos_\nu$ Satisfies Giraud Axioms]
\label{cor:noe-giraud}
$\InftyTopos_\nu$ satisfies all $\infty$-Giraud axioms.
\end{corollary}

%% ============================================================================
%% CHAPTER 5: NOETIC HOMOTOPY GROUPS
%% ============================================================================

\chapter{Noetic Homotopy Groups}
\label{ch:noetic-homotopy}

\begin{tcolorbox}[colback=blue!5!white,colframe=blue!75!black,title=\vnewfour]
This chapter develops noetic homotopy groups $\NoeticHomotopy$ and related
invariants.
\end{tcolorbox}

\section{Definition of Noetic Homotopy Groups}

\begin{definition}[Noetic Homotopy Groups]
\label{def:noetic-homotopy-groups}
For an Idea $\iota \in \Ideas$ and $n \geq 0$, the \textbf{$n$-th noetic homotopy group} is:
\[
\NoeticHomotopy(\iota) := \pi_n(\mathrm{Map}_{\mathcal{N}_\infty}(\iota, \iota), \mathrm{id}_\iota)
\]
\end{definition}

\begin{proposition}[Properties]
\label{prop:homotopy-properties}
\begin{enumerate}[label=(\roman*)]
  \item $\pi^\nu_0(\iota)$ classifies connected components of the Idea's automorphism space
  \item $\pi^\nu_1(\iota)$ is the fundamental group of automorphisms
  \item $\pi^\nu_n(\iota)$ for $n \geq 2$ are abelian
\end{enumerate}
\end{proposition}

\section{Noetic Whitehead Theorem}

\begin{theorem}[Noetic Whitehead Theorem]
\label{thm:noetic-whitehead}
A morphism $f : \iota \to \iota'$ in $\mathcal{N}_\infty$ is an equivalence iff
it induces isomorphisms on all noetic homotopy groups:
\[
f_* : \NoeticHomotopy(\iota) \xrightarrow{\cong} \NoeticHomotopy(\iota')
\]
for all $n \geq 0$.
\end{theorem}

\section{Postnikov Towers of Ideas}

\begin{definition}[Postnikov Tower]
\label{def:postnikov-tower}
The \textbf{Postnikov tower} of an Idea $\iota$ is a sequence:
\[
\cdots \to P_n(\iota) \to P_{n-1}(\iota) \to \cdots \to P_0(\iota)
\]
where $P_n(\iota)$ is the $n$-truncation with $\pi^\nu_k(P_n(\iota)) = 0$ for $k > n$.
\end{definition}

\begin{theorem}[Postnikov Convergence]
\label{thm:postnikov-conv}
\[
\iota \simeq \lim_n P_n(\iota)
\]
The Idea is recovered as the limit of its Postnikov tower.
\end{theorem}

\section{Homotopy Limits and Colimits}

\begin{definition}[Homotopy Limit]
\label{def:homotopy-limit}
The \textbf{homotopy limit} of a diagram $D : J \to \mathcal{N}_\infty$ is:
\[
\mathrm{holim}_J D := \mathrm{Map}(\mathcal{N}(J), \mathcal{N}_\infty) \times_{\mathrm{Map}(\partial\mathcal{N}(J), \mathcal{N}_\infty)} \{D\}
\]
\end{definition}

\begin{theorem}[Milnor Exact Sequence]
\label{thm:milnor}
For a tower of Ideas $\cdots \to \iota_2 \to \iota_1 \to \iota_0$:
\[
0 \to \lim^1 \pi^\nu_{n+1}(\iota_k) \to \pi^\nu_n(\mathrm{holim}_k \iota_k) \to \lim \pi^\nu_n(\iota_k) \to 0
\]
\end{theorem}

%% ============================================================================
%% CHAPTER 6: HIGHER STACKS AND MODULI
%% ============================================================================

\chapter{Higher Stacks and Moduli}
\label{ch:higher-stacks}

\begin{tcolorbox}[colback=blue!5!white,colframe=blue!75!black,title=\vnewfour]
This chapter develops noetic higher stacks and moduli problems.
\end{tcolorbox}

\section{Higher Stacks}

\begin{definition}[Higher Stack]
\label{def:higher-stack}
A \textbf{higher stack} (or $\infty$-stack) on the noetic site is a functor
\[
\mathcal{F} : \Ideas^{\mathrm{op}} \to \mathbf{Space}
\]
satisfying $\infty$-descent for all noetic coverings.
\end{definition}

\begin{definition}[Category of Noetic Stacks]
\label{def:stack-category}
The \textbf{category of noetic stacks} is:
\[
\HigherStack := \Sh_\infty(\Ideas, \tau_{\mathrm{noetic}})
\]
\end{definition}

\section{Moduli Problems for Ideas}

\begin{definition}[Moduli Stack of Ideas]
\label{def:moduli-ideas}
The \textbf{moduli stack of Ideas} $\mathcal{M}_{\Ideas}$ assigns to each Idea $\iota$
the space of Ideas ``over'' $\iota$:
\[
\mathcal{M}_{\Ideas}(\iota) := \mathrm{Map}_{\mathcal{N}_\infty}(\iota, \mathbf{Idea}_\infty)
\]
\end{definition}

\begin{definition}[Moduli of Fractals]
\label{def:moduli-fractals}
The \textbf{moduli stack of fractals} $\mathcal{M}_\Phi$ parametrizes fractal
structures on Ideas:
\[
\mathcal{M}_\Phi(\iota) := \{\Phi : \Phi \text{ is a fractal structure on } \iota\}
\]
\end{definition}

\section{Descent Conditions}

\begin{theorem}[ACBE Descent]
\label{thm:acbe-descent}
ACBE morphisms satisfy effective descent: for any ACBE covering
$\{U_\alpha \to X\}$, the natural functor
\[
\HigherStack(X) \to \mathrm{holim}_\Delta \HigherStack(U_\bullet)
\]
is an equivalence.
\end{theorem}

\section{Classifying Stacks}

\begin{definition}[Classifying Stack]
\label{def:classifying-stack}
The \textbf{classifying stack} $B\mathrm{Aut}(\iota)$ of an Idea $\iota$ represents
the functor:
\[
\iota' \mapsto \{\text{$\iota$-torsors over } \iota'\}
\]
\end{definition}

\begin{theorem}[Classification]
\label{thm:classification}
For connected $\iota'$:
\[
\mathrm{Map}(\iota', B\mathrm{Aut}(\iota)) \simeq \mathrm{Tors}_\iota(\iota')
\]
The mapping space classifies $\iota$-torsors.
\end{theorem}

%% ============================================================================
%% CHAPTER 7: SYNTHESIS AND CONNECTIONS
%% ============================================================================

\chapter{Synthesis and Connections}
\label{ch:synthesis}

\begin{tcolorbox}[colback=green!5!white,colframe=green!75!black,title=\textbf{Integration}]
This chapter connects the higher categorical framework to other v7.4 tracks.
\end{tcolorbox}

\section{Connection to Transfinite Fractals (Part V)}

\begin{proposition}[Fractal Functor to $\infty$-Categories]
\label{prop:fractal-infty}
The fractal functor $\FractalFunctor : \Ord \to \mathbf{Cat}$ extends to an
$\infty$-functor:
\[
\FractalFunctor_\infty : \Ord \to \InftyCat
\]
\end{proposition}

\section{Connection to Domain Semantics (Part VI)}

\begin{proposition}[Domain-$\infty$-Category Correspondence]
\label{prop:domain-infty}
Scott domains embed into $\infty$-groupoids via the nerve construction:
\[
\mathcal{N} : \ScottDom \hookrightarrow \InftyCat
\]
\end{proposition}

\section{Connection to v7.3 Meta Track}

\begin{theorem}[2-Category Upgrade]
\label{thm:2cat-upgrade}
The v7.3 2-categorical structures embed into $(\infty,1)$-categories:
\[
\TwoCat \hookrightarrow \InftyCat
\]
preserving adjunctions and limits.
\end{theorem}

\section{Summary Theorem}

\begin{theorem}[TKS Higher Categorical Foundations]
\label{thm:summary}
The noetic $(\infty,1)$-category $\mathbf{Idea}_\infty$ satisfies:
\begin{enumerate}[label=(\roman*)]
  \item It is an $\infty$-topos
  \item It admits all homotopy limits and colimits
  \item Noetic homotopy groups provide complete invariants
  \item Higher stacks classify moduli of noetic structures
  \item All v7.0--v7.3 structures embed faithfully
\end{enumerate}
\end{theorem}

\begin{tcolorbox}[colback=green!5!white,colframe=green!75!black,title=\textbf{Cross-Track Connections}]
\textbf{To Part V (Math)}: Ordinal-indexed $\infty$-categories via $\FractalFunctor_\infty$.

\textbf{To Part VI (Semantics)}: Domain-theoretic models as $\infty$-groupoids.

\textbf{To Part VIII (Geometry)}: Noetic manifolds as objects in a geometric
$\infty$-topos.
\end{tcolorbox}
